% English translation of capitulos/04b-qp2.tex, made by Claude Code (Claude Fable 5.2) on 24 November 2026.
% Pending review by the author. When the Portuguese file changes, update this one (scripts/verificar_traducao.py).
\subsection{RQ2: risks and mitigations}\label{sec:qp2}

This subsection synthesizes the risks and mitigations recorded in the extraction records of the 332 studies of the corpus. The records were extracted by AI and checked by the author in a sample of 67 studies; outside it, they remain without human checking, except for the claims verified one by one against the original. Each study enters once in each category it records. A mitigation is evaluated only when the study compares it with a condition without it, which does not imply effectiveness. The counts of the overview and of Table~\ref{tab:riscos} include all the records; those
of the following subsections come from the partial syntheses, made before the inclusion of S1610, S1617, S667 and of the studies of the last batch of stage 3.

\subsubsection{Overview}

% AFIRMACAO: A45
Functional correctness is recorded in 314 of the 332 studies, followed by cost (156), contamination (138), non-determinism (126), hallucination (99), security (90) and maintainability (41); 275 studies record at least one risk classified as ``other'' (Table~\ref{tab:riscos}).
% AFIRMACAO: A46
Almost all studies (328/332) record at least one mitigation, but only 144/332 evaluate any of them; per task, they are 112/247 in generation, 25/62 in debugging, 12/27 in translation and 7/20 in refactoring.

\begin{table}[!t]
\centering
\footnotesize
\setlength{\tabcolsep}{3pt}
\caption{Distinct studies per risk category and task}
\label{tab:riscos}
\begin{tabular}{p{2.6cm}p{0.75cm}p{0.75cm}p{0.75cm}p{0.75cm}p{0.75cm}}
\hline
Category & Gen. & Deb. & Tr. & Ref. & All \\
\hline
Studies (n) & 247 & 62 & 27 & 20 & 332 \\
\hline
Functional correctness & 232 & 58 & 27 & 18 & 314 \\
Security & 72 & 15 & 6 & 6 & 90 \\
Hallucination & 71 & 20 & 6 & 9 & 99 \\
Contamination & 96 & 31 & 13 & 6 & 138 \\
Non-determinism & 89 & 26 & 13 & 10 & 126 \\
Cost & 119 & 31 & 11 & 6 & 156 \\
Maintainability & 30 & 6 & 3 & 8 & 41 \\
Other & 202 & 54 & 22 & 16 & 275 \\
\hline
With mitigation & 244 & 61 & 27 & 20 & 328 \\
With evaluated mitigation & 112 & 25 & 12 & 7 & 144 \\
\hline
\end{tabular}
\par\smallskip
\parbox{\columnwidth}{\scriptsize Gen.: generation; Deb.: debugging; Tr.: translation; Ref.: refactoring. A study may have more than one task; the task columns do not add up to the total. Count by script from records extracted by AI, checked by sample. ``Evaluated mitigation'': at least one compared with a condition without it.}
\end{table}

\subsubsection{Synthesis by category}

\paragraph{Functional correctness}
% AFIRMACAO: A47
Passing tests is not equivalent to correctness: S1943 reports 62 correct patches among 93 plausible ones (S1943, Table I), S303 reports 186 plausible and 164 correct (S303, Table III) and, in S1562, tests generated by the agent passed wrong code in 90.9\% of the GPT failures (S1562, section 4.5).
% AFIRMACAO: A48
The prompting technique itself can worsen the measured outcome, with records for \textit{few-shot} (S6, section 4.2; S1598, p. 4, pass rate of 0.108 against 0.458 for \textit{zero-shot}), CoT in small models (S319, section 4.3; S1024, Table 5), RAG (S588, Fig. 5a), excess of context or of examples (S275, Figure 5; S1032, Fig. 3) and security preambles (S1693, section V-B).
% AFIRMACAO: A49
Self-correction and repair also make previously correct answers incorrect: S759 reports a $\Delta_{c \to i}$ of 8.4\% in Self-Refine with a small model (S759, Table 1), in S1592, the selection by reward model between the base generation and the hint-guided one corrupted 117 correct solutions and repaired 310 (S1592, section 3.3) and, in S2051, the rate after the first repair iteration stayed below the initial pass@3, measurements that are not equivalent, whereas selecting the best version improves the result (S2051, Table 8).

The mitigations evaluated were execution or compiler feedback loops (S302, S1053, S1077, S1086, S1562), external verification (S1977), interventions on the retrieved context (S670, S887) and control of reasoning according to difficulty (S1133, S1034), part of them by indirect indicators, such as similarity (S1155). Human review and manual checking of plausible patches remain as proposals (S1943, S2057).

\paragraph{Security}
The category brings together vulnerabilities in the generated code, privacy of the data sent to the model and operational security (S2060).
% AFIRMACAO: A50
When measured, vulnerabilities are frequent: in translation, from 28.6\% to 45.0\% of the translations of fixed sources introduce vulnerabilities and from 61.4\% to 71.9\% of the translations of vulnerable sources preserve them (S951, Table IV), and ``improvement'' iterations accumulated 387 distinct vulnerabilities, even with prompts that ask for more security (S739, Tables II and III).

The mitigations evaluated were security instructions in the prompt (S397, S427, S604, S951, S1179, S1195, S1604, S1607, S1693, S1763), injection of security knowledge (S695, S1894), reflection with detector feedback (S1065) and packages without isolation of mechanisms (S391, S2129).
% AFIRMACAO: A51
These evaluations record trade-offs with functional correctness: security guidelines reduced the Pass@1 of Claude-3.7-Sonnet from 16.19\% to 9.52\% (S1195, Table 7), RCI reduced pass@1 on HumanEval (S604, section V-C) and security preambles reduced correctness (S1693, section V-B); in S397, CWE-78 seemed not to be affected by the techniques (S397, section 6.3.2).

Sandbox or container (S137, S1415, S1694), human review between iterations (S739) and local deployment of the model (S1747, S1894) remain as proposals.

\paragraph{Hallucination}
The term has no uniform use; S511, for example, uses it for refactorings that change behaviour.
% AFIRMACAO: A52
A recurring pattern is the reference to non-existent APIs, functions, attributes or identifiers: in S1415, 47 to 48 of 100 reviewed failures are of this nature (S1415, Fig. 5), in S466, 22 cases, in a sample described as of 100 whose categories add up to 110 (S466, section 5.1), and in S109 at least one hardware value was hallucinated in 100\% of the IMU programs with GPT-3.5 (S109, section 5.2).
% AFIRMACAO: A53
Few evaluations measure hallucination itself: in S917, providing the list of signals in the prompt eliminated the non-existent signals observed with the base prompt, in a single task (S917, Tables V to X), whereas in S315 the inclusion of the imports was measured by pass@5 (0.233 without and 0.271 with the instruction, on 10 bugs) and in S1828, S1898 and S1977 the mitigations were measured by Pass@1 or pass rate.

\paragraph{Contamination}
% AFIRMACAO: A54
In the 128 studies that record it in the partial synthesis (138 in the current count, Table~\ref{tab:riscos}), contamination appears in general as a hypothesis; indications appear in S398 (at least 43 occurrences of each HumanEval prompt on GitHub, which indicates exposure, and not inclusion in the training of a specific model; section 3.1), S412 (24 of 98 correct patches lexically identical to the reference, section 6.1.1) and S2057 (122 of 522 bugs in the training data of StarCoder2, PDF p. 27).

The controls adopted (temporal cut-off, new benchmarks, renaming and similarity filters) are not compared with conditions without control.
% AFIRMACAO: A55
Only two studies evaluate a contamination mitigation against a condition without it: S137, by ablation with obfuscated prompts, whose outcome is performance and not contamination (Appendix M), and S2057, which compares trends between Defects4J and Defects4J-clean.

\paragraph{Non-determinism}
Repetitions, seeds and temperature 0 are adopted as a control, in general without comparison.
% AFIRMACAO: A56
Temperature 0 reduces, but does not eliminate, the variation between runs (S124, Table 5; S1695, Table V), and S782 reports a similarity of 0.7 between outputs of the same prompt with temperature 0 (S782, section 4.2); S124 evaluated two prompts: the \textit{concise prompt} made the candidates more deterministic than the base prompt at temperature 0, and CoT increased randomness in that condition (S124, Table 8).

\paragraph{Cost}
% AFIRMACAO: A57
Multi-step techniques add cost when it is measured: MapCoder uses on average 16.7 calls and 21.25 thousand tokens, against one call and 0.64 thousand tokens for the direct prompt (S344, Table 8), the technique of S2140 cost US\$ 0.421 against US\$ 0.017 for Default per 100 problems (S2140, section 7.1) and repair consumed from 5 to 15 times the energy of the single pass (S1415, section 6.1).

The units and the objects measured differ, which prevents aggregating the values; in some studies, the cost is that of the evaluation, and not of the prompting (S892, S1024).
% AFIRMACAO: A58
The cost mitigations evaluated use distinct metrics: candidate pruning reduced the time from 110.1 to 13.2 s per question (S195, Table 10), and routing by difficulty and \textit{adaptive} CoT were measured in tokens and in generation time, respectively (S1133, Table 2; S1034, Table 2).

\paragraph{Maintainability}
% AFIRMACAO: A59
When measured, maintainability gets worse in some scenarios: the C2C condition generated 62 \textit{blocker} or \textit{critical} code smells (56 and 6) against 17 (10 and 7) in the ground truth (S1708, Table VIII), complexity grew with the iterations in association with the vulnerabilities (r = 0.64; S739, section IV-D) and Zero-Verbosity reached a cyclomatic complexity of 54 (S1606, Table I).

Only two evaluated mitigations relate to the topic (S114, S2051).

\paragraph{Topics under ``other''}
In the partial syntheses, among the 248 studies that then recorded risks in the ``other'' category (275 in the current
count), threats to the validity of the studies themselves predominate: validity of the measurement and of the evaluation (100 of these
studies) and external validity (96). There are also context limits, unrequested behaviours (S1008) and human and ethical factors.
% AFIRMACAO: A60
In the latter, S62 reports participants who accepted without scrutiny the reasoning of incorrect patches (S62, section 5.4), and S140 records that \textit{few-shot} raised the social bias (CBS) of GPT-4 from 59.88\% to 68.56\% (S140, Table 7).

\subsubsection{Gaps}

\begin{itemize}
\item \textbf{Risks without evaluation.} Contamination (138 studies in the current count) has two mitigation evaluations in the
partial syntheses, and non-determinism (126), a single prompting one (S124). One possible reading is that the controls work as design care, without verifying the intended effect.
\item \textbf{Outcome different from the risk.} Hallucination mitigations are measured, in general, by Pass@k or pass rate; the term, like ``security'', covers distinct phenomena.
\item \textbf{Trade-offs between risks.} Few studies measure security and functional correctness together, and those that do record losses.
\item \textbf{Packages without isolation.} S1734, S2129 and S2140 do not allow the effect to be attributed to a component of the prompt.
\item \textbf{Uneven coverage.} Translation (27) and refactoring (20) have few studies.
\item \textbf{Checking by sample.} Counts and locators depend on records that, outside the sample of 67 studies, were not checked by a human.
\end{itemize}
